% Auto-generated complete chapter from 28D_artstation_portfolio_benchmark (1).md
\chapter{28D — ArtStation / Portfolio Benchmark}
\label{complete:root_034}

\begin{sourcemeta}
\textbf{Fuente:} \href{file:///E:/Laboral/28D_artstation_portfolio_benchmark\%20\%281\%29.md}{28D\_artstation\_portfolio\_benchmark (1).md}\\
\textbf{Seccion:} root\\
\textbf{Estado:} duplicate\\
\textbf{Rol:} Version parcial / duplicada\\
\textbf{Lineas originales:} 271\\
\textbf{Ultima modificacion:} 2026-06-11 18:06\\
\textbf{Deduplicacion conservadora:} 0 bloques repetidos omitidos; 0 caracteres repetidos no reimpresos.
\end{sourcemeta}

\section*{Resumen de orientacion}
--- tipo: modulo\_benchmark modulo: 28D\_artstation\_portfolio\_benchmark estado: corregido fecha: 2026-06-11 fuente\_base: A4\_artstation\_portfolio\_benchmark perfil\_objetivo: Alexander Woodcock — Real-Time 3D / Unity WebGL / Technical Visualization revision: 2 nota: reemplaza la versión anterior porque A4 no investigó ArtStation de forma válida ---

\section*{Contenido completo depurado}
\addcontentsline{toc}{section}{Contenido completo depurado}

\medskip\hrule\medskip


tipo: modulo\_benchmark

modulo: 28D\_artstation\_portfolio\_benchmark

estado: corregido

fecha: 2026-06-11

fuente\_base: A4\_artstation\_portfolio\_benchmark

perfil\_objetivo: Alexander Woodcock — Real-Time 3D / Unity WebGL / Technical Visualization

revision: 2

nota: reemplaza la versión anterior porque A4 no investigó ArtStation de forma válida

\medskip\hrule\medskip







\subsection{1. Corrección crítica}




La versión anterior de este módulo no era válida como \textbf{ArtStation benchmark}.




El archivo A4 usado como base contenía portafolios externos, blogs técnicos y GitHub Pages, pero no una muestra real de perfiles o proyectos de ArtStation. Por tanto, cualquier conclusión presentada como patrón de ArtStation queda retirada.




Este módulo corregido separa tres cosas:




\begin{enumerate}
\item lo que sí sirve del intento anterior;
\item lo que no debe usarse;
\item la tarea necesaria para completar ArtStation correctamente.
\end{enumerate}




\subsection{2. Estado de evidencia}




\begin{center}
\scriptsize
\begin{longtable}{>{\raggedright\arraybackslash}p{0.305\linewidth} >{\raggedright\arraybackslash}p{0.305\linewidth} >{\raggedright\arraybackslash}p{0.305\linewidth}}
\toprule
Bloque & Evidencia & Estado \\
\midrule
\endfirsthead
\toprule
Bloque & Evidencia & Estado \\
\midrule
\endhead
ArtStation perfiles & 0 & inválido \\
ArtStation proyectos & 0–1 & insuficiente \\
Blogs técnicos & 3+ & usable \\
Portfolios externos & 7+ & usable \\
GitHub Pages & 4+ & usable \\
Benchmark ArtStation & no & pendiente \\
Benchmark portfolio general & parcial & usable \\
\bottomrule
\end{longtable}
\end{center}




\subsection{3. Lo que queda invalidado}




No deben usarse como conclusiones de ArtStation:




\begin{itemize}
\item “los perfiles de ArtStation usan X estructura”;
\item “los breakdowns de ArtStation priorizan X”;
\item “ArtStation demuestra X patrón de contratación”;
\item “Alexander debe adaptar X desde ArtStation”.
\end{itemize}




La muestra no sostiene esas afirmaciones.




\subsection{4. Lo que sí puede conservarse}




El intento anterior sí aporta patrones generales de portfolio técnico.




\begin{center}
\scriptsize
\begin{longtable}{>{\raggedright\arraybackslash}p{0.225\linewidth} >{\raggedright\arraybackslash}p{0.225\linewidth} >{\raggedright\arraybackslash}p{0.225\linewidth} >{\raggedright\arraybackslash}p{0.225\linewidth}}
\toprule
Fuente & Plataforma & Uso válido & Fuerza \\
\midrule
\endfirsthead
\toprule
Fuente & Plataforma & Uso válido & Fuerza \\
\midrule
\endhead
Potion Shader & Blog & shader breakdown & alta \\
Harry Alisavakis & Blog & autoridad técnica & alta \\
julhe & GitHub Pages & posts shader & alta \\
Francisco Múrias & GitHub Pages & grid profesional & media \\
Alson Entuna & GitHub Pages & rol especializado & media \\
Alla Eddine Rakik & Framer & portfolio técnico & media \\
Omar Rinaz & GitHub Pages & CV visible & baja \\
\bottomrule
\end{longtable}
\end{center}




\subsection{5. Patrones generales válidos}




\subsubsection{5.1. Resultado visual primero}




Un caso técnico debe abrir con el resultado:




\begin{mdcode}
Bloque de codigo: text
final render / gif / video / interactive demo / before-after
\end{mdcode}




Para Alexander, esto aplica a:




\begin{itemize}
\item TwinSight X500;
\item cross-section shader;
\item exploded view;
\item Inspect mode;
\item Analyze mode;
\item thermal heuristic view;
\item CAD-to-WebGL pipeline.
\end{itemize}




\subsubsection{5.2. Rol explícito}




El portfolio debe decir qué hizo Alexander.




Formato recomendado:




\begin{mdcode}
Bloque de codigo: text
Role: Solo developer / technical artist / 3D pipeline owner
\end{mdcode}




No basta con mostrar el proyecto como tesis.




\subsubsection{5.3. Breakdown por problema}




Cada case study debe responder:




\begin{mdcode}
Bloque de codigo: text
Problem → Constraint → Process → Technical solution → Result → Limitation
\end{mdcode}




Esto es más fuerte que una galería visual.




\subsubsection{5.4. Herramientas visibles}




El stack debe aparecer en el hero o primera sección:




\begin{mdcode}
Bloque de codigo: text
Unity · WebGL · Blender · C# · UI Toolkit · Shader Graph/HLSL · 3D Optimization
\end{mdcode}




\subsubsection{5.5. Limitaciones declaradas}




Para perfiles técnicos, declarar límites aumenta credibilidad.




Ejemplos aplicables:




\begin{itemize}
\item thermal view no es FEA;
\item visual twin no es digital twin operativo;
\item validación formativa si la muestra es pequeña;
\item WebGL implica restricciones de memoria y carga.
\end{itemize}




\subsection{6. Estructura recomendada para TwinSight X500}




\begin{mdcode}
Bloque de codigo: text
1. Hero visual
2. One-line summary
3. Role / scope / tools
4. Problem
5. Constraints
6. CAD-to-WebGL pipeline
7. Interaction systems
8. Shader / visualization systems
9. Optimization notes
10. Validation evidence
11. Limitations
12. Links
\end{mdcode}




\subsection{7. Case studies derivados}




\begin{center}
\scriptsize
\begin{longtable}{>{\raggedright\arraybackslash}p{0.305\linewidth} >{\raggedright\arraybackslash}p{0.305\linewidth} >{\raggedright\arraybackslash}p{0.305\linewidth}}
\toprule
Case study & Función & Prioridad \\
\midrule
\endfirsthead
\toprule
Case study & Función & Prioridad \\
\midrule
\endhead
TwinSight X500 & flagship & 1 \\
CAD-to-WebGL pipeline & pipeline & 2 \\
Cross-section shader & shader & 3 \\
Exploded view system & interaction & 4 \\
Bottom-sheet UI & UX/system & 5 \\
Fastener inspection & detail system & 6 \\
Thermal heuristic view & analysis view & 7 \\
\bottomrule
\end{longtable}
\end{center}




\subsection{8. Estructura de página recomendada}




\begin{mdcode}
Bloque de codigo: markdown
# TwinSight X500

## Summary
Interactive WebGL technical visualization of a Holybro X500 V2 drone assembly.

## Role
Solo developer / technical artist / 3D pipeline owner.

## Tools
Unity, WebGL, Blender, C#, UI Toolkit, Shader Graph/HLSL.

## Problem
2D manuals require users to mentally reconstruct spatial relationships.

## Constraints
Browser deployment, mobile-first UI, real-time interaction, academic validation.

## Technical Breakdown
- CAD cleanup and hierarchy normalization.
- Runtime part selection.
- Exploded view.
- Cross-section clipping.
- Bottom-sheet technical data.
- Heuristic thermal visualization.

## Result
A public WebGL prototype for inspection and study of a drone assembly.

## Limitations
Visual product twin, not operational digital twin.
Thermal mode is heuristic, not calibrated FEA.

## Links
Demo, repository, final report, manuals.
\end{mdcode}




\subsection{9. Qué evitar}




\begin{center}
\scriptsize
\begin{longtable}{>{\raggedright\arraybackslash}p{0.455\linewidth} >{\raggedright\arraybackslash}p{0.455\linewidth}}
\toprule
Error & Riesgo \\
\midrule
\endfirsthead
\toprule
Error & Riesgo \\
\midrule
\endhead
Grid sin breakdown & parece galería \\
Proyecto como “tesis” solamente & reduce valor laboral \\
Reclamar digital twin operativo & sobreventa \\
Mostrar métricas no cerradas & riesgo técnico \\
Ocultar limitaciones & baja credibilidad \\
No incluir video/gif & baja evidencia visual \\
\bottomrule
\end{longtable}
\end{center}




\subsection{10. Decisión operativa}




Antes de convertir este módulo en benchmark final, ejecutar:




\begin{mdcode}
Bloque de codigo: text
A4_retry_artstation_only_benchmark
\end{mdcode}




Hasta entonces, 28D debe usarse como:




\begin{itemize}
\item guía de portfolio técnico general;
\item auditoría de error;
\item estructura provisional de case studies;
\item no como benchmark ArtStation final.
\end{itemize}




\subsection{11. Prompt para A4-Retry}




\begin{mdcode}
Bloque de codigo: text
Execute A4_retry_artstation_only_benchmark.

Context:
The previous A4 attempt failed. It collected portfolio/blog examples but did not properly research ArtStation profiles or ArtStation project breakdowns.

Objective:
Collect 20–35 public ArtStation examples relevant to Alexander’s technical art and real-time 3D positioning.

Prioritize:
- Technical Artist profiles;
- Unity Technical Artist profiles;
- shader breakdown projects;
- VFX breakdown projects;
- real-time 3D breakdowns;
- optimization breakdowns;
- Blender technical breakdowns;
- XR / interactive 3D projects;
- technical visualization projects;
- portfolios with tools, GitHub, demo videos or detailed process notes.

Output structured tables only.

Fields:
- Artist / profile label
- ArtStation URL
- Source type: profile / resume / project / blog post / channel result
- Role target
- Seniority if visible
- Project type
- Breakdown structure
- Images used
- Videos used
- Technical text depth
- Tags / skills
- Tools shown
- Portfolio / GitHub links if visible
- What Alexander can adapt
- What Alexander should not copy
- Confidence
- Date checked

Rules:
- ArtStation sources only.
- Use public pages only.
- Do not invent missing data.
- If unavailable, write “not found”.
- If ArtStation blocks access, record “blocked” and do not infer details.
- Include source URLs.
- Tables only.
- No narrative strategy.
\end{mdcode}



